\section{Parameters and Integration}
\label{sec:parameters}

This section explains how parameters interact across the ARC-SPRUCE stack: the lattice signature (trapdoor and
preimages), the rational hash, the SmallWood proof system, and the PRF. Detailed tables and extended benchmarks are
deferred to Appendix~\ref{app:ext-params}.

\subsection{Parameter dependencies}
All components share the same base ring \(\Rq\) and modulus \(q\):
\begin{itemize}[leftmargin=1.25em]
  \item The signature map \(A\in\Rq^{n\times m}\) and hash key \(B\) are defined over \(\Rq\).
  \item SmallWood constraints are evaluated over the base field \(\Fq\) (coefficients modulo \(q\)) and therefore must
  use the same \(q\).
  \item The PRF is instantiated over \(\Fq\) so that its constraints are native to the same field.
\end{itemize}
The coefficient bound \(\BoundMsg\) governs:
(i) sampling ranges for \(\mu,\veck,r_0^U,r_1^U\) and issuer challenges,
(ii) the centering modulus \(\Delta=2\BoundMsg+1\), and
(iii) membership polynomials used in the NIZKs.

\subsection{Sampler parameters}
The trapdoor sampler must output short signatures \(u\) with overwhelming probability and with a distribution close to
a discrete Gaussian over the appropriate lattice coset. This requires choosing:
\begin{itemize}[leftmargin=1.25em]
  \item ring degree \(N\) and modulus \(q\),
  \item trapdoor quality (e.g., an NTRU trapdoor parameter \(\alpha\)),
  \item per-slot Gaussian widths \(\sigma\) (or \(\sigma_k[i]\) for a two-plane sampler),
  \item an acceptance predicate (norm bound) that yields a public verification bound \(\BoundSig\).
\end{itemize}
We use an Antrag-based NTRU trapdoor and a hybrid two-plane Gaussian sampler; details and tuning rules are given in
Appendix~\ref{app:gaussian-sampler}.

\subsection{vSIS / rational hash parameters}
The rational hash must satisfy:
\begin{itemize}
  \item \textbf{Denominator invertibility.} The denominator \(B_3-r_1\) must be invertible component-wise; the holder
  enforces this by rejection sampling \(r_1\) (and the NIZK proves boundedness and the cleared relation).
  \item \textbf{No wrap-around in bounds.} Membership checks for \([-\BoundMsg,\BoundMsg]\) are implemented via vanishing
  polynomials over \(\Fq\), so we require \(\BoundMsg \ll q\) to avoid modular wrap-around artifacts.
  \item \textbf{Compatibility with signature hardness.} The vSIS analyses \cite{EPRINT:DKLW25,PKC:DKLW25} impose
  conditions on the rational hash family (e.g., strong linear independence and denominator bounds) which translate into
  constraints on the distribution of bounded randomness and on the public key structure.
\end{itemize}

\subsection{SmallWood parameters}
SmallWood introduces parameters controlling packing, masking, and soundness:
\begin{itemize}[leftmargin=1.25em]
  \item \textbf{Packing size} \(s=|\Omega|\) (often chosen as \(s=n_{\mathrm{cols}}\)).
  \item \textbf{Tail length} \(\ell\) for HVZK masking and degree enforcement.
  \item \textbf{Spot-check count} \(\ell'\) and number of masks \(\rho\) for PIOP soundness.
  \item \textbf{Degree bounds} \(d_{\mathrm{decs}}=s+\ell-1\) for witness rows and \(d_Q\) for constraint-derived masks.
  \item \textbf{Small-field knob} \(\theta\) when using an extension-field variant.
\end{itemize}
The pre-sign proof is dominated by linear constraints (commitment, centering, and a linear cleared-denominator hash),
while the post-sign proof introduces higher effective degree due to the PRF’s S-box constraints and the bilinear hash
relation (since \(T\) becomes a witness row). Appendix~\ref{app:smallwood} and Appendix~\ref{app:ext-params} detail how
these degrees translate into \(d_Q\) and final soundness.

\subsection{PRF parameters}
The PRF must be:
(i) secure at target level \(\lambda\),
(ii) efficiently representable in constraints, and
(iii) compatible with \(\Fq\).
We use a Poseidon2-like permutation with feed-forward and truncation as a concrete instantiation. Its parameters
include width \(t=\ell_{\mathsf{key}}+\ell_{\nonce}\), round counts \((R_F,R_P)\), MDS matrices, round constants,
S-box exponent \(d\), and truncation length \(\ell_{\prftag}\). We require \(\ell_{\prftag}\ge
\lceil \lambda/\log_2(q)\rceil\) so the tag contains \(\lambda\) bits of entropy. The algorithm is given in
Section~\ref{subsec:poseidon-prf}, while the checkpointed constraint shape used in the post-sign proof is stated in
Section~\ref{subsec:prf-checkpointed}. Appendix~\ref{app:prf} records extended notes (optional).

\subsection{Sizes and performance (example instantiation)}
To ground the parameter interaction, we record one concrete demo instantiation (ring degree \(N=1024\), modulus
\(q=1038337\), target \(\lambda\approx 128\) bits) and report measured proof sizes for the two NIZKs. The pre-sign proof
uses only the base witness rows, while the post-sign proof adds signature and PRF rows as described in
Section~\ref{subsec:smallwood-showing}.

\begin{table}[H]
\centering
\renewcommand{\arraystretch}{1.1}
\begin{tabular}{@{}lccc@{}}
\toprule
\multirow{2}{*}{\textbf{Proof}} & \textbf{Rows} & \multirow{2}{*}{\textbf{Size}} & \textbf{Soundness} \\
 & \textbf{(base / extra)} &  & \textbf{(bits)} \\
\midrule
Pre-sign NIZK (issue) & \(9\) / \(0\) & \(\approx 7.4\) KB & \(\approx 129.6\) \\
Post-sign NIZK (show) & \(12\) / \(86\) & \(\approx 20.0\) KB & \(\approx 129.4\) \\
\bottomrule
\end{tabular}
\caption{Example proof sizes and row counts for one instantiation; see Appendix~\ref{app:ext-params} for parameter tuples
and extended data.}
\label{tab:proof-sizes}
\end{table}
